% 基本信息--------------------------------------------------------------------------------
% 学位类型在 main.tex 的 \documentclass[...]{zufe} 中指定：
%   bachelor（默认，本科毕业论文/专业实践报告）、master（硕士）、doctor（博士）

% 【仅本科生效】指定论文类型 0:毕业论文  1/2:专业实践1/2
% 硕士、博士模板不使用该变量，无需填写
\ifzufeBachelor
  \newcommand{\reportStyle}{0}
\fi

% 在这里填写你的论文中文题目
\newcommand{\thesisTitle}{海岛环境对武学宗师成长的影响机理}
% 若摘要页题目与封面不同，则取消下一行注释，并填写摘要页使用的中文题目
% 由于封面标题和摘要标题的字体和内容差异，这个需求通常发生在封面标题需要自定义合理断行，
% 但摘要页不需要断行的情况。注：断行使用双反斜杠\\
% 在这里定义摘要页题目（如果与封面标题相同，则注释掉下一行）
% \newcommand{\thesisTitleAbs}{海岛环境对武学宗师成长的影响机理}
% 在这里填写你的论文英文题目
\newcommand{\thesisTitleEN}{Influence mechanism of island environment on the growth of martial arts masters}

% 若有副标题，则运行下一行的代码，若无副标题，则将下一行注释掉（在\haveSub{}最前面添加 % 号）
\haveSub{}

% 在这里填写你的论文中文副标题（没有只需注释掉\haveSub{}即可）
\newcommand{\thesisSubTitle}{基于桃花岛武学流派的研究}
% 在这里填写你的论文英文副标题（没有只需注释掉\haveSub{}即可）
\newcommand{\thesisSubTitleEN}{Research Based on Taohua island martial arts school}

% 在这里填写你的相关信息
% 本科封面用到：作者姓名、学号、指导教师、所在学院、专业名称、班级
% 硕士/博士封面用到：作者姓名、专业、所在学院、指导教师
\newcommand{\deptName}{信息技术与人工智能学院}
\newcommand{\majorName}{软件工程}
\newcommand{\yourName}{xx}
\newcommand{\yourStudentID}{xxxxxxxxxxxx}
\newcommand{\mentorName}{xxx}
\newcommand{\className}{xxx}
% 完成日期（中文）：用于封面、中文扉页
\newcommand{\Today}{20xx年6月}
% 完成日期（英文）：用于英文扉页
\newcommand{\TodayEN}{June 20xx}
% 基本信息--------------------------------------------------------------------------------

% 中英文摘要与关键词
\newcommand{\abstractCN}{
  本文……。\textcolor{blue}{摘要的内容要包括研究的目的、方法、结果和结论。计量单位一律换算成国际标准计量单位。除特殊情况外，数字一律用阿拉伯数字。中、英文摘要的内容应严格一致。}
}

\newcommand{\keywordsCN}{
  计算机科学与技术；算法；复杂度；深度学习；机器学习；可视化计算
}

\newcommand{\abstractEN}{
  In order to study...... The content should include the purpose, methods, results and conclusions of the research. All measurement units should be converted to international standard units. Unless otherwise specified, all numbers should be expressed in Arabic numerals. The contents of the “Chinese” and ``English'' abstracts should be strictly consistent.
}

\newcommand{\keywordsEN}{
  Computer science and technology; Algorithm; Complexity; Deep learning; Machine learning; Visualization analysis
}